\documentclass[addpoints]{exam}
% leqno 將方程式編號放在左側

%\printanswers

\usepackage[top=2.5cm,bottom=2cm,left=2.5cm,right=2.5cm,headsep=10pt,a4paper]{geometry} % 引用頁面幾何套件，設定頁面邊界

\usepackage{siunitx} % 角度
% \input{NewCommands}

\usepackage{amsmath,amssymb,amsthm} % 引用常見的數學符號與設定


\theoremstyle{definition}
\newtheorem*{definition}{Definition}
\newtheorem*{theorem}{Theorem}
\newtheorem*{remark}{Remark}
\newcommand{\R}{\mathbb{R}}

\usepackage{lastpage} % 取得最後頁碼




\usepackage{graphicx} % 引用圖檔內嵌入套件
\graphicspath{{Pictures/}} % 設定圖檔目錄
\usepackage{xcolor} %引用顏色套件
\usepackage{enumerate}
\usepackage{float}
\usepackage{hyperref}
\usepackage{mdwlist} % use suspend & resume to successive enumerate

\usepackage{CJKutf8} % 設定中文

\def \lflr{\left\lfloor} %自定義 floor function
\def \rflr{\right\rfloor} %自定義

% question separation
\renewcommand{\questionshook}{%
  \setlength{\itemsep}{0.5cm}%
}

\allowdisplaybreaks %equation allowed to next page

\firstpageheader{\today}
                {}
                {}
\footer{}{\sf Page \thepage~of~\pageref{LastPage}}{}

\begin{document}

\begin{CJK*}{UTF8}{gbsn}
% Title
\begin{center}
    {\Large{\bf Introduction to Mathematical Analysis II\\
    Homework 7  Due  May 1st (Friday), 2026\\
    Please submit your group homework online in PDF format.\\
You are expected to work collaboratively within your group. 
Please do not obtain complete solutions directly from AI tools. 
The purpose of this assignment is to develop your own mathematical reasoning and problem-solving skills. %  Homework X Brief Solution
  }} \\ 
    \large
    %\text{Due date: 3/9/2023}
\end{center}

\noindent\rule{16.2cm}{0.4pt}


\begin{questions}


% ============ Question x ============
\question  (15 pts)  (Exercises 7.5.5.)
Let $f:\mathbf{R}^n\to\mathbf{R}$ be Lebesgue measurable, and let $g:\mathbf{R}^n\to\mathbf{R}$
be a function which agrees with $f$ outside of a set of measure zero, thus there exists a set
$A\subseteq\mathbf{R}^n$ of measure zero such that $f(x)=g(x)$ for all $x\in\mathbf{R}^n\setminus A$.
Show that $g$ is also Lebesgue measurable.
(\emph{Hint:} use Exercise~7.4.10.)

\begin{solution}
\end{solution}

\question (15 pts) 
Suppose that \(f:\mathbf{R}\to [0,\infty)\) is Lebesgue measurable and
\(g:[0,\infty)\to [0,\infty)\) is monotone. Prove that
\(g\circ f\) is Lebesgue measurable.



\begin{solution}


\end{solution}



\question (20 pts) A function \(f:\mathbf{R}\to \mathbf{R}\) is \emph{upper semicontinuous} if
\[
\lim_{k\to\infty} x_k=x
\quad\Longrightarrow\quad
\limsup_{k\to\infty} f(x_k)\le f(x).
\]


\begin{enumerate}
    \item[(a)] Draw a graph of an upper semicontinuous function that is not continuous.
    
    \item[(b)] Show that upper semicontinuity is equivalent to the requirement that for every open ray \((-\infty,a)\), the preimage
    \[
    f^{-1}((-\infty,a))
    \]
    is an open set.
     \item[(c)] Show that an upper semicontinuous function is measurable.
    
\end{enumerate}


\begin{solution} 

\end{solution}

\question (15 pts)  Let $\Omega$ be a measurable subset of $\mathbb{R}^n$, and let $f:\Omega\to[0,\infty]$ and $g:\Omega\to[0,\infty]$ be measurable functions.
Without using Theorem 8.2.9 or Lemma 8.2.10, prove that
$$
\int_\Omega (f+g)\ge \int_\Omega f + \int_\Omega g.
$$

\begin{solution} 
 
 \end{solution}
 
 
\question (15) Let \(\{s_n\}\) be an increasing sequence of nonnegative simple functions which converges pointwise on an interval \(I\) to a limit function \(f\). If \(I\) is unbounded and if \(f(x)\ge 1\) almost everywhere on \(I\), prove that the sequence
\[
\left\{\int_I s_n\right\}
\] diverges.

\begin{solution} 
 
 \end{solution}



\question (20 pts)
Let \(\{s_n\}_{n=1}^\infty\) be a sequence of measurable simple functions on \(\mathbf{R}\) such that
\[
s_n(x)\to f(x)\qquad \text{for every }x\in \mathbf{R}.
\]
\begin{enumerate}
    \item[(a)] Prove that for every \(a\in \mathbf{R}\),
    \[
    f^{-1}\bigl((a,\infty)\bigr)
    =
    \bigcup_{n=1}^\infty \bigcap_{k=n}^\infty
    s_k^{-1}\!\left(\left(a+\frac1n,\infty\right)\right).
    \]
    
    \item[(b)] Deduce that \(f\) is measurable.
\end{enumerate}


   
\begin{solution} 
 
 \end{solution}



\end{questions}
\end{CJK*}
\end{document}

